\documentclass[../../../include/open-logic-section]{subfiles}

\begin{document}

\olfileid[hi]{sth}{choice}{justifications}
\olsection{चयन के बारे में आंतरिक विचार}

तब व्यापक प्रश्न यह है कि सुव्यवस्था, अथवा चयन, अथवा वास्तव में आकार की
दृष्टि से सभी समुच्चयों की तुलनीयता---इनमें से किसे लें, इससे अंतर नहीं
पड़ता---का औचित्य दिया जा सकता है या नहीं।

यहाँ एक \emph{आंतरिक} औचित्य देने का प्रयास है। \olref[z][story]{sec} में
हमने पदानुक्रम के बारे में कई सिद्धांत प्रस्तुत किए थे। उनमें से एक को फिर
कहना उपयोगी है:
\begin{enumerate}
	\item[] \stagesacc। किसी भी चरण $S$ और चरण $S$ से \emph{पहले} बने
	किन्हीं समुच्चयों के लिए: चरण $S$ पर ऐसा समुच्चय बनता है जिसके सदस्य ठीक
	वे समुच्चय हैं। चरण $S$ पर और कुछ नहीं बनता।
\end{enumerate}
वास्तव में अनेक लेखकों ने सुझाया है कि चयन अभिगृहीत का औचित्य इस सिद्धांत
(जैसे किसी सिद्धांत) के माध्यम से दिया जा सकता है। हम उस दृष्टिकोण की एक
संक्षिप्त व्याख्या देंगे।

हम एक छोटे-से सरल परिणाम से आरंभ करेंगे, जो चयन का \emph{एक और} तुल्य रूप
देता है:

\begin{thm}[$\ZF$ में]\ollabel{choiceset}
चयन निम्न सिद्धांत के तुल्य है। यदि $A$ के !!{element} असंयुक्त और अरिक्त
हैं, तो कोई $C$ है ऐसा कि प्रत्येक $x\in A$ के लिए $C\cap x$ एकल समुच्चय
है। (ऐसे $C$ को हम $A$ के लिए {चयन समुच्चय} कहते हैं।)
\end{thm}

इस परिणाम का प्रमाण सीधा है और हम इसे पाठक के अभ्यास के लिए छोड़ते हैं।

\begin{prob}
\olref[sth][choice][justifications]{choiceset} सिद्ध कीजिए। यदि कठिनाई हो,
तो \cite[pp.~242--3]{Potter2004} में प्रमाण मिल सकता है।
\end{prob}

मुख्य बात यह है कि $A$ के लिए चयन समुच्चय, $A$ के लिए किसी चयन फलन का
परास मात्र है। अतः चयन का औचित्य देने के लिए हम चयन समुच्चयों के अस्तित्व
के पदों में उसके तुल्य सूत्रीकरण का औचित्य देने का प्रयास कर सकते हैं। अब हम
ठीक यही करेंगे।

$A$ के !!{element} असंयुक्त और अरिक्त हों। \stageshier{}
(\olref[z][story]{sec} देखें) से $A$ किसी चरण $S$ पर बनता है। ध्यान दें कि
$\bigcup A$ के सभी !!{element} चरण $S$ से पहले उपलब्ध होते हैं। अब
\stagesacc{} से $S$ से पहले बने \emph{किन्हीं भी} समुच्चयों के लिए ऐसा
समुच्चय बनता है जिसके सदस्य ठीक वे समुच्चय हैं। दूसरे शब्दों में: पहले
उपलब्ध समुच्चयों का प्रत्येक \emph{संभव} संग्रह चरण $S$ पर अस्तित्व में
होगा। किंतु ऐसी वस्तुओं का चयन करना निश्चय ही \emph{संभव} है जिनसे $A$ के
लिए चयन समुच्चय बन सके; वह $\bigcup A$ का कोई अत्यंत विशिष्ट उपसमुच्चय
मात्र है। अतः आवश्यकतानुसार ऐसा कोई चयन समुच्चय अस्तित्व में है।

यह आंतरिक आधार पर चयन का औचित्य देने का \emph{बहुत} संक्षिप्त प्रयास है।
किंतु इस विचार को आगे बढ़ाने के लिए आपको पॉटर का इसका सुंदर विकास
(\citeyear[\S14.8]{Potter2004}) पढ़ना चाहिए।

\end{document}
